\documentclass[11pt]{article}
\usepackage[margin=1in]{geometry}
\usepackage{amsmath,amssymb,amsthm,mathtools}
\usepackage{enumitem}
\usepackage{microtype}
\input{glyphtounicode}
\pdfgentounicode=1
\pdfglyphtounicode{mu}{03BC}
\pdfglyphtounicode{bardbl}{2016}
\pdfglyphtounicode{parenleftbig}{0028}
\pdfglyphtounicode{parenrightbig}{0029}
\pdfglyphtounicode{parenleftbigg}{0028}
\pdfglyphtounicode{parenrightbigg}{0029}
\pdfglyphtounicode{parenlefttp}{239B}
\pdfglyphtounicode{parenleftbt}{239D}
\pdfglyphtounicode{parenrighttp}{239E}
\pdfglyphtounicode{parenrightbt}{23A0}
\pdfglyphtounicode{circleplusdisplay}{2A01}
\pdfglyphtounicode{summationdisplay}{2211}
\pdfglyphtounicode{productdisplay}{220F}
\usepackage[colorlinks=true,linkcolor=blue,citecolor=blue,urlcolor=blue]{hyperref}
\usepackage{accsupp}
\NewCommandCopy{\OriginalHookrightarrow}{\hookrightarrow}
\NewCommandCopy{\OriginalLongmapsto}{\longmapsto}
\NewCommandCopy{\OriginalUpbracefill}{\upbracefill}
\renewcommand{\hookrightarrow}{\mathrel{
  \BeginAccSupp{method=hex,unicode,ActualText=21AA}
  \OriginalHookrightarrow\EndAccSupp{}}}
\renewcommand{\longmapsto}{\mathrel{
  \BeginAccSupp{method=hex,unicode,ActualText=27FC}
  \OriginalLongmapsto\EndAccSupp{}}}
\renewcommand{\upbracefill}{
  \BeginAccSupp{method=hex,unicode,ActualText=23DF}
  \OriginalUpbracefill\EndAccSupp{}}
\hypersetup{pdftitle={Quadratic stabilization of the canonical Foulkes--Howe map},pdfauthor={Rewritten for readability from the OpenAI manuscript of the same title}}
\newtheorem{theorem}{Theorem}
\newtheorem{corollary}[theorem]{Corollary}
\newtheorem{proposition}[theorem]{Proposition}
\newtheorem{lemma}[theorem]{Lemma}
\theoremstyle{definition}
\newtheorem{convention}[theorem]{Convention}
\theoremstyle{remark}
\newtheorem{remark}[theorem]{Remark}
\DeclareMathOperator{\Sym}{Sym}
\DeclareMathOperator{\Av}{Av}
\newcommand{\C}{\mathbb C}
\newcommand{\Sa}{\mathfrak S_a}
\newcommand{\D}[2]{D_{#1,#2}}
\newcommand{\cP}{\mathcal P}
\newcommand{\unverified}[1]{\par\smallskip\noindent\fbox{\parbox{0.95\linewidth}{\textbf{Unverified step.} #1}}\par\smallskip}
\title{Quadratic stabilization of the canonical Foulkes--Howe map}
\author{Rewritten for readability from the OpenAI manuscript of the same title}
\date{}
\begin{document}
\maketitle
\begin{abstract}
For every finite-dimensional complex vector space $V$, the canonical
Foulkes--Howe map $\Sym^b(\Sym^a V)\longrightarrow\Sym^a(\Sym^b V)$ is
surjective whenever $a\ge2$ and $b\ge a(a-1)$. This is a quadratic
stabilization bound independent of $\dim V$.
\end{abstract}

\section{The result and how it is proved}\label{sec:introduction}

\subsection{Setting}

Throughout, $V$ is a finite-dimensional complex vector space and
$\Sym V=\bigoplus_{k\ge0}\Sym^kV$ is its symmetric algebra. A product of
vectors always means their product in $\Sym V$. An element
$u_1\cdots u_k\in\Sym^kV$ with $u_i\in V$ is called a \emph{product of $k$
vectors} (the vectors may repeat or be zero), and $u^k$ is a \emph{pure
power}.

\paragraph{Averaging.}
Let $W$ be a complex vector space and $a\ge1$ an integer. The symmetric group
$\Sa$ acts on $W^{\otimes a}$ by permuting tensor factors. Define the
linear map $\Av\colon W^{\otimes a}\to W^{\otimes a}$ by
\begin{equation}\label{eq:average}
 \Av(w_1\otimes\cdots\otimes w_a)
 =\frac1{a!}\sum_{\pi\in\Sa}
 w_{\pi(1)}\otimes\cdots\otimes w_{\pi(a)}.
\end{equation}
We use three facts about it.
\begin{enumerate}[label=(\roman*),nosep]
\item $\Av\tau$ is $\Sa$-invariant for every $\tau$, and $\Av\tau=\tau$ if
$\tau$ is invariant (all $a!$ permuted copies of $\tau$ equal $\tau$).
Hence $\Av$ is a projection onto $(W^{\otimes a})^{\Sa}$, and the averages
of simple tensors span $(W^{\otimes a})^{\Sa}$.
\item The right side of \eqref{eq:average} is symmetric in
$(w_1,\ldots,w_a)$, so $\Av$ factors through the quotient map
$\rho\colon W^{\otimes a}\to\Sym^aW$: there is a linear map
$\overline{\Av}\colon\Sym^aW\to(W^{\otimes a})^{\Sa}$ with
$\overline{\Av}(w_1\cdots w_a)=\Av(w_1\otimes\cdots\otimes w_a)$.
\item $\overline{\Av}$ is an isomorphism with inverse
$\rho|_{(W^{\otimes a})^{\Sa}}$. Indeed,
$\rho(\Av(w_1\otimes\cdots\otimes w_a))=\frac1{a!}\sum_\pi
w_{\pi(1)}\cdots w_{\pi(a)}=w_1\cdots w_a$, and for invariant $\tau$,
$\overline{\Av}(\rho(\tau))=\Av(\tau)=\tau$ by (i).
\end{enumerate}
We call $\overline{\Av}$ the \emph{averaging isomorphism}. Note that
$\overline{\Av}(w^a)=w^{\otimes a}$.

\paragraph{The algebra $R$ and the map.}
Fix an integer $a\ge1$. Let $\cP=(\Sym V)^{\otimes a}$ with componentwise
multiplication
$(f_1\otimes\cdots\otimes f_a)(g_1\otimes\cdots\otimes g_a)
=f_1g_1\otimes\cdots\otimes f_ag_a$. It is a commutative algebra on which
$\Sa$ acts by algebra automorphisms. (The source calls this algebra $P_a$;
it is renamed here so that $P$ can denote a product of vectors.) For $j\ge0$
put
\[
 R_j=\bigl((\Sym^j V)^{\otimes a}\bigr)^{\Sa}\subseteq\cP,\qquad
 R=\bigoplus_{j\ge0}R_j .
\]
A product of invariants is invariant and
$\Sym^iV\cdot\Sym^jV\subseteq\Sym^{i+j}V$, so $R_iR_j\subseteq R_{i+j}$.
Thus $R$ is a graded commutative algebra. Let $A\subseteq R$ be the
subalgebra generated by $R_1$ and put $A_b=A\cap R_b$. Since a product of
$j$ elements of $R_1$ lies in $R_j$, the space $A_b$ is the span of all
products of exactly $b$ elements of $R_1$.

The averaging isomorphism with $W=V$ and with $W=\Sym^bV$ gives
$R_1\cong\Sym^aV$ and $R_b\cong\Sym^a(\Sym^bV)$. For $b\ge0$, the
multiplication $(z_1,\ldots,z_b)\mapsto z_1\cdots z_b$ is a symmetric
multilinear map $R_1^{\,b}\to R_b$, so it induces a linear map
$\Sym^bR_1\to R_b$ with image $A_b$. Transported through the two averaging
isomorphisms, this is the \emph{canonical Foulkes--Howe map}
\begin{equation}\label{eq:map}
 \mu_{a,b,V}:\Sym^b(\Sym^a V)\longrightarrow\Sym^a(\Sym^b V),
 \qquad \operatorname{im}\mu_{a,b,V}=A_b .
\end{equation}
Concretely, $u^a$ corresponds to $u^{\otimes a}\in R_1$ and
$(u_1^{\otimes a})\cdots(u_b^{\otimes a})=(u_1\cdots u_b)^{\otimes a}$, so
$\mu_{a,b,V}\bigl((u_1^a)\cdots(u_b^a)\bigr)=(u_1\cdots u_b)^a$.

\begin{theorem}\label{thm:main}
For every integer $a\ge2$, every integer $b\ge a(a-1)$, and every
finite-dimensional complex vector space $V$, the map
$\mu_{a,b,V}$ in \eqref{eq:map} is surjective. Equivalently, $A_b=R_b$.
\end{theorem}

For $a=1$ we have $R=\Sym V$, and $\mu_{1,b,V}$ is the identity map in every
degree.

\subsection{Proof at a glance}

\paragraph{Key idea.}
By duality (Proposition~\ref{prop:criterion}), surjectivity of
$\mu_{a,b,V}$ is equivalent to the following statement: \emph{a symmetric
$a$-linear form $T$ on $\Sym^bV$ with $T(P,\ldots,P)=0$ for every product
$P$ of $b$ vectors is zero.} This is proved by induction on $a$
(Proposition~\ref{prop:vanishing}). The tool is
Lemma~\ref{lem:shift}: the directional-derivative operator
$\D YU=\sum_kY_k\,\partial/\partial U_k$ is injective on polynomials whose
degree in the block $U$ exceeds their degree in the block $Y$. The lemma is
used backwards. To show that a polynomial is zero, we apply a chain of such
operators until we reach a polynomial already known to be zero, and we check
the degree inequality before each operator.

\paragraph{The inductive step.}
Put $r=a-1$ and $m=b-r$; the hypothesis $b\ge a(a-1)$ says $m\ge r^2$. Let
$Q$ be a product of $m$ vectors and $x,t\in V$.
\begin{enumerate}[label=\textbf{Step \arabic*.},leftmargin=*,nosep]
\item Over $\C$, $x^r+ct^r$ is a product of $r$ vectors, so
$(x^r+ct^r)Q$ is a product of $b$ vectors for every $c\in\C$. The
coefficient of $c$ in the diagonal value gives
$T(x^rQ,\ldots,x^rQ,t^rQ)=0$.
\item A chain of $m$ operators, each with a different target block $y_i$,
turns $T(x^rQ,\ldots,x^rQ,t^b)$ into a multiple of the expression in Step~1.
Lemma~\ref{lem:shift} gives $T(x^rQ,\ldots,x^rQ,t^b)=0$.
\item Step~2 is the diagonal hypothesis for the symmetric $r$-linear form
$(Q_1,\ldots,Q_r)\mapsto T(x^rQ_1,\ldots,x^rQ_r,t^b)$ on $\Sym^mV$. The
induction hypothesis (which needs $m\ge r(r-1)$ and $m\ge1$) shows that this
form is zero.
\item A chain of $r^2$ operators, all with the same target block $x$, turns
$T(v_1^b,\ldots,v_r^b,t^b)$ into a multiple of
$T(x^rv_1^m,\ldots,x^rv_r^m,t^b)$, which is zero by Step~3. Before each
operator the $x$-degree is at most $r^2-1$ and the degree of the
differentiated block is at least $m+1\ge r^2+1$, so
Lemma~\ref{lem:shift} gives $T(v_1^b,\ldots,v_r^b,t^b)=0$. Pure powers span
$\Sym^bV$, hence $T=0$.
\end{enumerate}
In Step~4 the degree of the single shared target $x$ accumulates to $r^2$;
this is where the quadratic bound enters.

\paragraph{Smallest case $a=2$.}
Here $r=1$, $b\ge2$ and $m=b-1$. Step~1 gives $T(xQ,tQ)=0$. Step~2 gives
$T(xQ,t^b)=0$ for products $Q$ of $b-1$ vectors. Step~3 is the case $a=1$
(products span $\Sym^{b-1}V$) and gives $T(xQ',t^b)=0$ for all
$Q'\in\Sym^{b-1}V$. Step~4 applies the single operator $\D x{v_1}$ to
$T(v_1^b,t^b)$, producing $b\,T(xv_1^{b-1},t^b)=0$ with degrees $b>0$.

\paragraph{Dependencies.}
\begin{itemize}[nosep]
\item Lemma~\ref{lem:span} (pure powers span) $\Rightarrow$
Proposition~\ref{prop:criterion} and the base case and last line of
Proposition~\ref{prop:vanishing}.
\item Lemma~\ref{lem:poly}, Convention~\ref{conv:poly} and
Lemma~\ref{lem:rule} (bookkeeping for polynomials built from $T$)
$\Rightarrow$ Steps~2 and~4.
\item Lemma~\ref{lem:shift} (differential injectivity) $\Rightarrow$
Steps~2 and~4.
\item Proposition~\ref{prop:criterion} $+$ Proposition~\ref{prop:vanishing}
$\Rightarrow$ Theorem~\ref{thm:main} $\Rightarrow$
Corollary~\ref{cor:comparison}.
\end{itemize}
Section~\ref{sec:context} collects the source's remarks on the literature;
nothing there is used in the proofs.

\section{Preliminaries}\label{sec:prelim}

\begin{lemma}[Pure powers span]\label{lem:span}
For every integer $k\ge1$, the space $\Sym^kV$ is spanned by products of
$k$ vectors, and also by pure powers $u^k$, $u\in V$.
\end{lemma}

\noindent\emph{Used in:} Proposition~\ref{prop:criterion} (with $k=a$), and
the base case and final line of Proposition~\ref{prop:vanishing} (with
$k=b$).

\begin{proof}
Products span because $\Sym^kV$ is the image of $V^{\otimes k}$, and simple
tensors map to products. For pure powers we prove the polarization identity
\begin{equation}\label{eq:polarization}
 k!\,u_1\cdots u_k
 =\sum_{J\subseteq\{1,\ldots,k\}}(-1)^{k-|J|}
   \Bigl(\sum_{j\in J}u_j\Bigr)^k
 \qquad(u_1,\ldots,u_k\in V).
\end{equation}
First take formal commuting variables $s_1,\ldots,s_k$ in place of the
$u_j$. By the multinomial theorem, $(\sum_{j\in J}s_j)^k$ is the sum of
$\binom{k}{\beta}s^\beta$ over exponent vectors $\beta$ with $|\beta|=k$ and
support in $J$. Let $\beta$ have support $S$. Its coefficient on the right
of \eqref{eq:polarization} is
\[
 \binom k\beta\sum_{J\supseteq S}(-1)^{k-|J|}
 =\binom k\beta(-1)^{k-|S|}\sum_{J'\subseteq\{1,\ldots,k\}\setminus S}(-1)^{|J'|}
 =\binom k\beta(-1)^{k-|S|}(1-1)^{k-|S|}.
\]
This is zero unless $S=\{1,\ldots,k\}$. In that case $\beta=(1,\ldots,1)$
and the coefficient is $k!$. So \eqref{eq:polarization} holds in
$\C[s_1,\ldots,s_k]$. Applying the algebra homomorphism
$\C[s_1,\ldots,s_k]\to\Sym V$, $s_j\mapsto u_j$, gives it in $\Sym V$. Since
$k!\ne0$ in $\C$, every product of $k$ vectors is a linear combination of
pure powers.
\end{proof}

\begin{lemma}[Two facts about complex polynomials]\label{lem:poly}
Let $F\in\C[Z_1,\ldots,Z_N]$.
\begin{enumerate}[label=(\alph*),nosep]
\item If $F$ vanishes at every point of $\C^N$, then $F=0$.
\item Let $U$ be a subset of the variables and $d\ge0$. If
$F(\lambda U,\text{rest})=\lambda^dF(U,\text{rest})$ at every point and for
every $\lambda\in\C$, then $F$ is homogeneous of degree $d$ in $U$.
\end{enumerate}
\end{lemma}

\begin{proof}
(a) Induct on $N$. Write $F=\sum_jF_j(Z_1,\ldots,Z_{N-1})Z_N^j$. For fixed
values of $Z_1,\ldots,Z_{N-1}$ this is a univariate polynomial with
infinitely many roots, so every $F_j$ vanishes at that point. By induction
$F_j=0$.

(b) Write $F=\sum_eF_e$ with $F_e$ homogeneous of degree $e$ in $U$. At
each point, $\sum_e(\lambda^e-\lambda^d)F_e=0$ for all $\lambda$, so
$F_e$ vanishes at that point for $e\ne d$. By (a), $F_e=0$ for $e\ne d$.
\end{proof}

\begin{convention}[Vector variables and polynomials built from $T$]
\label{conv:poly}
Assume $n=\dim V>0$ and fix a basis $e_1,\ldots,e_n$ of $V$. A \emph{vector
variable} $u$ is a block of $n$ independent scalar variables
$(u_1,\ldots,u_n)$, standing for the vector $\sum_ku_ke_k$. We use the same
letter for the block and for the vector. Let $T$ be an $a$-linear form on
$\Sym^bV$, and consider an expression
\[
 F=T(M_1,\ldots,M_a),
\]
in which each $M_s$ is a product of $b$ vectors, each factor being a fixed
vector or one of finitely many vector variables. Expanding every vector
variable in the basis and using multilinearity of the product and of $T$
shows that $F$ is a polynomial function of the scalar coordinates. By
Lemma~\ref{lem:poly}(a) it is a uniquely determined polynomial, and
``$F=0$ for all vector values'' is the same as ``$F=0$ as a polynomial''.
If a vector variable $u$ occurs $d$ times in total among
$M_1,\ldots,M_a$ (counted with multiplicity), then replacing $u$ by
$\lambda u$ multiplies $F$ by $\lambda^d$, so by
Lemma~\ref{lem:poly}(b) $F$ is homogeneous of degree $d$ in the block $u$.

For two vector variables $U,Y$, define the operator
\[
 \D{Y}{U}=\sum_{k=1}^n Y_k\frac{\partial}{\partial U_k}.
\]
It lowers the $U$-degree by one and raises the $Y$-degree by one.
\end{convention}

\begin{lemma}[Differentiation rule]\label{lem:rule}
In the situation of Convention~\ref{conv:poly}, let $u$ and $y$ be distinct
vector variables. Suppose $u$ occurs in exactly one slot $s$, where
$M_s=u^jz$ with $j\ge1$ and $z$ a product of $b-j$ vectors not involving
$u$. Then
\[
 \D yu\,T(M_1,\ldots,M_a)
 =j\,T(M_1,\ldots,M_{s-1},\,u^{j-1}y\,z,\,M_{s+1},\ldots,M_a).
\]
\end{lemma}

\noindent\emph{Used in:} Steps~2 and~4 of
Proposition~\ref{prop:vanishing}.

\begin{proof}
Fix values of all vector variables and let $\phi(\epsilon)$ be the value of
$F=T(M_1,\ldots,M_a)$ after replacing $u$ by $u+\epsilon y$ (and nothing
else). By the chain rule, $\phi'(0)=\sum_ky_k\,\partial F/\partial u_k
=(\D yuF)$ at the chosen point. On the other hand, the binomial theorem in
the commutative algebra $\Sym V$ gives
$(u+\epsilon y)^jz=u^jz+j\epsilon\,u^{j-1}yz+\epsilon^2w(\epsilon)$ with
$w(\epsilon)\in\Sym^bV$ polynomial in $\epsilon$. Only slot $s$ changes
and $T$ is linear in it, so
$\phi(\epsilon)=F+j\epsilon\,T(\ldots,u^{j-1}yz,\ldots)+\epsilon^2(\cdots)$.
Hence $\phi'(0)=j\,T(\ldots,u^{j-1}yz,\ldots)$. The two sides of the claim
agree at every point, hence as polynomials by Lemma~\ref{lem:poly}(a).
\end{proof}

\section{The canonical image and diagonal vanishing}\label{sec:map}

This section turns surjectivity of $\mu_{a,b,V}$ into a vanishing statement
about symmetric multilinear forms. It is used once, in the proof of
Theorem~\ref{thm:main}.

\begin{proposition}[Diagonal criterion]\label{prop:criterion}
Let $a,b\ge1$ be integers. The map $\mu_{a,b,V}$ is surjective if and only
if every symmetric $a$-linear form $T$ on $\Sym^b V$ satisfying
$T(P,\ldots,P)=0$ for all products $P$ of $b$ vectors is zero.
\end{proposition}

\begin{proof}
Put $E=\Sym^bV$, so $R_b=(E^{\otimes a})^{\Sa}$.

\emph{Step A: the elements $z(u):=u^{\otimes a}$, $u\in V$, span $R_1$.}
The averaging isomorphism $\Sym^aV\to R_1$ sends $u^a$ to $u^{\otimes a}$,
and the pure powers $u^a$ span $\Sym^aV$ by Lemma~\ref{lem:span} with
$k=a$.

\emph{Step B: $A_b=\operatorname{span}_{\C}\{P^{\otimes a}:
P\text{ a product of }b\text{ vectors}\}$.}
By definition $A_b$ is spanned by the products $z_1\cdots z_b$ with
$z_i\in R_1$. The product is multilinear, so by Step~A it is enough to take
$z_i=z(u_i)$. Componentwise multiplication gives
$z(u_1)\cdots z(u_b)=(u_1\cdots u_b)^{\otimes a}$.

\emph{Step C: functionals on $R_b$ correspond bijectively to symmetric
$a$-linear forms on $E$.}
For $\lambda\in R_b^*$ define
\[
 T_\lambda(w_1,\ldots,w_a)
 =\lambda\bigl(\Av(w_1\otimes\cdots\otimes w_a)\bigr)
 \qquad(w_i\in E).
\]
This is a symmetric $a$-linear form. The map $\lambda\mapsto T_\lambda$ is
injective because averages of simple tensors span $R_b$ (fact~(i) on
averaging). It is surjective: given a symmetric $a$-linear form $T$, let
$\widetilde T\colon E^{\otimes a}\to\C$ be its linearization and
$\lambda=\widetilde T|_{R_b}$. Then
$T_\lambda(w_1,\ldots,w_a)=\frac1{a!}\sum_\pi
T(w_{\pi(1)},\ldots,w_{\pi(a)})=T(w_1,\ldots,w_a)$ by symmetry of $T$.

\emph{Step D: $T_\lambda(P,\ldots,P)=\lambda(P^{\otimes a})$ for every
$P\in E$.} Indeed $\Av(P^{\otimes a})=P^{\otimes a}$.

\emph{Step E: conclusion.}
By Steps~B and~D, $\lambda$ annihilates $A_b$ if and only if
$T_\lambda(P,\ldots,P)=0$ for every product $P$ of $b$ vectors. By Step~C,
the annihilator of $A_b$ in $R_b^*$ is zero if and only if every symmetric
form with vanishing diagonal on products is zero. A subspace $A_b$ of the
finite-dimensional space $R_b$ equals $R_b$ if and only if its annihilator
is zero, and $A_b=\operatorname{im}\mu_{a,b,V}$ by \eqref{eq:map}.
\end{proof}

\section{Differential injectivity}\label{sec:differential}

The following lemma is the only analytic input. The source notes that it is
the positive-weight injectivity in the usual $\mathfrak{sl}_2$ weight
strings; compare
\cite[Section~11.1, Claim~11.4 and equation~(11.5)]{FultonHarris}. The proof
below is self-contained apart from the definition of the Fischer inner
product. In the source the adjoint operator is called $E$; it is renamed
$D'$ here, because $E$ denotes $\Sym^bV$ in Section~\ref{sec:map}.

\begin{lemma}[Differential injectivity]\label{lem:shift}
Let $n\ge1$, let $U=(U_1,\ldots,U_n)$ and $Y=(Y_1,\ldots,Y_n)$ be two
disjoint blocks of variables, and let $p,q$ be nonnegative integers. On
polynomials homogeneous of degree $p$ in $U$ and degree $q$ in $Y$, the
operator $\D{Y}{U}$ is injective if $p>q$. The polynomials may depend on any
finite collection of additional variables independent of $U$ and $Y$.
\end{lemma}

\noindent\emph{Used in:} Steps~2 and~4 of
Proposition~\ref{prop:vanishing}, each time on a polynomial of the stated
bidegree, with all other blocks as additional variables.

\begin{proof}
\emph{Plan.} We put an inner product on the polynomial ring for which the
adjoint of $D:=\D YU$ is $D':=\D UY$. We compute the commutator $D'D-DD'$,
and read off a norm identity that forces $h=0$ when $Dh=0$ and $p>q$.

\emph{Step 1: the inner product.}
List all scalar variables present (those of $U$, of $Y$, and the additional
ones) as $Z_1,\ldots,Z_N$. Give $\C[Z_1,\ldots,Z_N]$ the Fischer, or apolar,
Hermitian inner product
\cite[Section~2, equation~(8) and Proposition~3]{AldazRender}:
\[
 \Bigl\langle\sum_\alpha a_\alpha Z^\alpha,\sum_\beta b_\beta Z^\beta\Bigr\rangle
 =\sum_\alpha\alpha!\,a_\alpha\overline{b_\alpha},
 \qquad\alpha!:=\prod_k\alpha_k!.
\]
Distinct monomials are orthogonal, $\|Z^\alpha\|^2=\alpha!>0$, and the form
is positive definite.

\emph{Step 2: differentiation is adjoint to multiplication.}
For every $i$ and all polynomials $g,h$,
$\langle\partial_{Z_i}g,h\rangle=\langle g,Z_ih\rangle$. Both sides are
linear in $g$ and conjugate-linear in $h$, so it suffices to take
$g=Z^\alpha$, $h=Z^\beta$. Writing $e_i$ for the $i$-th unit exponent
vector, both sides equal $\alpha!$ if $\alpha=\beta+e_i$ and $0$ otherwise
(on the left, $\alpha_i\,(\alpha-e_i)!=\alpha!$). Taking complex conjugates
gives also $\langle Z_ig,h\rangle=\langle g,\partial_{Z_i}h\rangle$.

\emph{Step 3: $\langle Dg,h\rangle=\langle g,D'h\rangle$.}
By Step~2, used twice,
\[
 \langle Dg,h\rangle=\sum_k\langle Y_k\,\partial_{U_k}g,h\rangle
 =\sum_k\langle\partial_{U_k}g,\partial_{Y_k}h\rangle
 =\sum_k\langle g,U_k\,\partial_{Y_k}h\rangle=\langle g,D'h\rangle .
\]

\emph{Step 4: the commutator.}
By the product rule,
\begin{align*}
 D'D&=\sum_{j,k}U_j\partial_{Y_j}Y_k\partial_{U_k}
 =\sum_kU_k\partial_{U_k}+\sum_{j,k}U_jY_k\partial_{Y_j}\partial_{U_k},\\
 DD'&=\sum_{j,k}Y_k\partial_{U_k}U_j\partial_{Y_j}
 =\sum_kY_k\partial_{Y_k}+\sum_{j,k}Y_kU_j\partial_{U_k}\partial_{Y_j}.
\end{align*}
Subtracting,
\[
 D'D-DD'=N_U-N_Y,
 \qquad
 N_U=\sum_k U_k\partial_{U_k},\quad
 N_Y=\sum_k Y_k\partial_{Y_k}.
\]

\emph{Step 5: the norm identity.}
Let $h$ be homogeneous of degree $p$ in $U$ and $q$ in $Y$. On each of its
monomials $N_U$ acts as multiplication by $p$ and $N_Y$ by $q$, so
$(N_U-N_Y)h=(p-q)h$. By Step~3,
$\|Dh\|^2=\langle D'Dh,h\rangle$ (the number is real) and
$\|D'h\|^2=\langle DD'h,h\rangle$. Hence, by Step~4,
\begin{equation}\label{eq:norm}
 \|Dh\|^2-\|D'h\|^2=\langle(N_U-N_Y)h,h\rangle=(p-q)\|h\|^2.
\end{equation}

\emph{Step 6: conclusion.}
If $Dh=0$ and $p>q$, the left side of \eqref{eq:norm} is $-\|D'h\|^2\le0$
and the right side is $\ge0$. So $(p-q)\|h\|^2=0$, and $h=0$ by positive
definiteness.
\end{proof}

\section{Vanishing on products}\label{sec:vanishing}

\begin{proposition}\label{prop:vanishing}
Let $V$ be a finite-dimensional complex vector space, and let $a\ge1$ and
$b\ge1$ be integers satisfying $b\ge a(a-1)$.
If a symmetric $a$-linear form $T$ on $\Sym^b V$ satisfies
\begin{equation}\label{eq:diagonal}
 T(P,\ldots,P)=0
 \qquad\text{for every product $P$ of $b$ vectors},
\end{equation}
then $T=0$.
\end{proposition}

\begin{proof}
\emph{Plan.} We induct on $a$, proving the statement for all admissible $b$
and all $V$ at once. In the inductive step, Step~1 extracts one identity
from \eqref{eq:diagonal}, Steps~2 and~4 each use Lemma~\ref{lem:shift} along
a chain of derivatives, and Step~3 applies the induction hypothesis in
between.

If $V=0$, then $\Sym^bV=0$ because $b\ge1$, and $T=0$. So assume
$\dim V>0$ and fix a basis, so that Convention~\ref{conv:poly} applies.

\paragraph{Base case $a=1$.}
$T$ is a linear form on $\Sym^bV$ vanishing on all products of $b$ vectors.
These span $\Sym^bV$ (Lemma~\ref{lem:span}), so $T=0$.

\paragraph{Inductive step.}
Let $a\ge2$, assume the proposition for $a-1$ (all $b$, all $V$), and put
\begin{equation}\label{eq:rm}
 r=a-1\ge1,\qquad m=b-r.
\end{equation}
From $b\ge a(a-1)=r(r+1)$ we get
\begin{equation}\label{eq:range}
 m=b-r\ge r^2\ge1 .
\end{equation}

\paragraph{Step 1: a split binary form.}
\emph{Claim:} for all $x,t\in V$ and every product $Q=y_1\cdots y_m$ of $m$
vectors,
\begin{equation}\label{eq:binary}
 T(\underbrace{x^rQ,\ldots,x^rQ}_{r\text{ slots}},t^rQ)=0.
\end{equation}

(1a) \emph{For every $c\in\C$, $P_c:=(x^r+ct^r)Q$ is a product of $b$
vectors.} In $\C[\xi,\eta]$ (the source writes $X,Y$ for these two scalar
variables), if $c\ne0$ choose $\alpha\in\C$ with $\alpha^r=-c$ and put
$\zeta=e^{2\pi i/r}$. Since
$\prod_{j=0}^{r-1}(\xi-\zeta^j\omega)=\xi^r-\omega^r$, substituting
$\omega=\alpha\eta$ gives
\[
 \xi^r+c\eta^r=\prod_{j=0}^{r-1}(\xi-\alpha\zeta^j\eta).
\]
For $c=0$ the factorization is $\xi^r$. Applying the algebra homomorphism
$\C[\xi,\eta]\to\Sym V$, $\xi\mapsto x$, $\eta\mapsto t$ (no independence of
$x$ and $t$ is needed) shows that $x^r+ct^r$ is a product of $r$ vectors of
$V$. Multiplying by $Q$ gives a product of $r+m=b$ vectors.

(1b) \emph{Extracting the coefficient of $c$.} By (1a) and
\eqref{eq:diagonal}, $T(P_c,\ldots,P_c)=0$ for all $c\in\C$. Write
$P_c=x^rQ+c\,t^rQ$. By multilinearity and symmetry,
\[
 T(P_c,\ldots,P_c)=\sum_{k=0}^a\binom akc^k\,
 T(\underbrace{x^rQ,\ldots,x^rQ}_{a-k},\underbrace{t^rQ,\ldots,t^rQ}_{k}).
\]
A polynomial in $c$ vanishing for all $c\in\C$ has zero coefficients. The
coefficient of $c^1$ is $a\,T(x^rQ,\ldots,x^rQ,t^rQ)$ and $a\ne0$, which
gives \eqref{eq:binary}.

\paragraph{Step 2: freeing the last slot.}
\emph{Claim:} for all $x,t\in V$ and every product $Q=y_1\cdots y_m$ of $m$
vectors,
\begin{equation}\label{eq:step2}
 H_x:=T(\underbrace{x^rQ,\ldots,x^rQ}_{r\text{ slots}},t^b)=0.
\end{equation}

Fix $x\in V$ and regard $t,y_1,\ldots,y_m$ as independent vector variables.
For $0\le k\le m$ put
\[
 G_k=T(x^rQ,\ldots,x^rQ,\;t^{b-k}y_1\cdots y_k),
\]
so $G_0=H_x$. By Convention~\ref{conv:poly}, $G_k$ is a polynomial,
homogeneous of degree $b-k$ in $t$, of degree $r+1$ in $y_i$ for $i\le k$,
and of degree $r$ in $y_i$ for $i>k$.

(2a) \emph{$G_m=0$.} Since $b-m=r$, we have
$G_m=T(x^rQ,\ldots,x^rQ,t^rQ)$. This is zero for all vector values by
\eqref{eq:binary}, hence zero as a polynomial.

(2b) \emph{$\D{y_{k+1}}{t}G_k=(b-k)\,G_{k+1}$ for $0\le k<m$.} The variable
$t$ occurs only in the last slot, as $t^{b-k}\cdot(y_1\cdots y_k)$ with
$b-k\ge r+1\ge1$. Lemma~\ref{lem:rule} applies with $u=t$, $y=y_{k+1}$. In
particular
\begin{equation}\label{eq:first-composition}
 \D{y_k}{t}\cdots\D{y_1}{t}H_x
 =\frac{b!}{(b-k)!}\,
 T(x^rQ,\ldots,x^rQ,t^{b-k}y_1\cdots y_k)\qquad(0\le k\le m).
\end{equation}

(2c) \emph{Descending induction: $G_{k+1}=0$ implies $G_k=0$.} If
$G_{k+1}=0$ then $\D{y_{k+1}}{t}G_k=0$ by (2b). Apply
Lemma~\ref{lem:shift} with $U=t$, $Y=y_{k+1}$ and the other $y_i$ as
additional variables. The polynomial $G_k$ has
\[
 p=\deg_t=b-k,\qquad q=\deg_{y_{k+1}}=r,\qquad p-q=m-k\ge1.
\]
Here $y_{k+1}$ still has degree $r$, because the earlier operators raised
only the degrees of $y_1,\ldots,y_k$. Hence $G_k=0$.

Starting from (2a) and applying (2c) for $k=m-1,\ldots,0$ gives
$G_0=H_x=0$ as a polynomial, hence for all vector values.

\paragraph{Step 3: induction on the remaining slots.}
\emph{Claim:} for all $x,t\in V$,
\begin{equation}\label{eq:inductive-vanishing}
 T(x^rQ_1,\ldots,x^rQ_r,t^b)=0
 \qquad(Q_1,\ldots,Q_r\in\Sym^m V).
\end{equation}
Fix $x,t\in V$ and define
$S_{x,t}(Q_1,\ldots,Q_r):=T(x^rQ_1,\ldots,x^rQ_r,t^b)$. Multiplication by
$x^r$ is a linear map $\Sym^mV\to\Sym^bV$ (as $r+m=b$), so $S_{x,t}$ is a
symmetric $r$-linear form on $\Sym^mV$. By \eqref{eq:step2}, its diagonal
vanishes at every product $Q$ of $m$ vectors. By \eqref{eq:range},
$m\ge1$ and $m\ge r^2\ge r(r-1)$, so the induction hypothesis (the
proposition for $a-1=r$, with $m$ in place of $b$ and the same $V$) gives
$S_{x,t}=0$.

\paragraph{Step 4: removing the common power.}
\emph{Claim:} $T(v_1^b,\ldots,v_r^b,t^b)=0$ for all
$v_1,\ldots,v_r,t\in V$.

Fix $t\in V$ and regard $v_1,\ldots,v_r$ and $x$ as independent vector
variables. Until now $x$ was a fixed vector; this is legitimate because
\eqref{eq:inductive-vanishing} holds for every $x\in V$. For a tuple
$\mathbf j=(j_1,\ldots,j_r)$ with $0\le j_i\le r$ put
\[
 F_{\mathbf j}=T\bigl(x^{j_1}v_1^{b-j_1},\ldots,x^{j_r}v_r^{b-j_r},\;t^b\bigr).
\]
By Convention~\ref{conv:poly}, $F_{\mathbf j}$ is homogeneous of degree
$b-j_i$ in $v_i$ and of degree $j_1+\cdots+j_r$ in $x$. The polynomial of
interest is $f:=F_{(0,\ldots,0)}=T(v_1^b,\ldots,v_r^b,t^b)$, of $x$-degree
zero.

(4a) \emph{$F_{(r,\ldots,r)}=0$.} Since $b-r=m$, this is
$T(x^rv_1^m,\ldots,x^rv_r^m,t^b)$. It is zero for all vector values by
\eqref{eq:inductive-vanishing} with $Q_i=v_i^m$, hence as a polynomial.

(4b) \emph{$\D x{v_i}F_{\mathbf j}=(b-j_i)\,F_{\mathbf j+\mathbf e_i}$
whenever $j_i\le r-1$,} where $\mathbf e_i$ is the $i$-th unit tuple. The
variable $v_i$ occurs only in slot $i$, as $v_i^{b-j_i}\cdot x^{j_i}$ with
$b-j_i\ge m+1\ge1$. Lemma~\ref{lem:rule} applies with $u=v_i$, $y=x$. The
derivative does not touch the factors $x$ already present or any other
slot.

(4c) \emph{The chain.} Raise $j_1$ from $0$ to $r$ one unit at a time,
then $j_2$, and so on up to $j_r$. This is a chain of $r^2$ tuples steps
from $(0,\ldots,0)$ to $(r,\ldots,r)$, and it corresponds to applying
$\D x{v_1}^{\,r}$, then $\D x{v_2}^{\,r}$, \ldots, then $\D x{v_r}^{\,r}$.
Multiplying the constants in (4b) gives
\begin{equation}\label{eq:second-composition}
 \D{x}{v_r}^{\,r}\cdots\D{x}{v_1}^{\,r}f
 =\left(\frac{b!}{m!}\right)^r
 T(x^rv_1^m,\ldots,x^rv_r^m,t^b)=0,
\end{equation}
because each slot contributes $b(b-1)\cdots(b-r+1)=b!/m!$.

(4d) \emph{Descending induction along the chain: $F_{\mathbf j+\mathbf
e_i}=0$ implies $F_{\mathbf j}=0$} for each step
$\mathbf j\to\mathbf j+\mathbf e_i$ of the chain. If
$F_{\mathbf j+\mathbf e_i}=0$ then $\D x{v_i}F_{\mathbf j}=0$ by (4b). Apply
Lemma~\ref{lem:shift} with $U=v_i$, $Y=x$ and the other blocks as additional
variables. Before this operator, $v_i$ has been differentiated
$j_i\le r-1$ times and at most $r^2-1$ operators have been applied in
total, so by \eqref{eq:range}
\[
 p=\deg_{v_i}F_{\mathbf j}=b-j_i\ge b-r+1=m+1\ge r^2+1,
 \qquad
 q=\deg_xF_{\mathbf j}=j_1+\cdots+j_r\le r^2-1 .
\]
Thus $p>q$ and $F_{\mathbf j}=0$.

Starting from (4a) and descending the chain gives $f=0$. As $t$ was
arbitrary, the claim of Step~4 follows.

\paragraph{Conclusion of the inductive step.}
By Step~4, $T(w_1,\ldots,w_a)=0$ whenever every $w_i$ is a pure power
$u_i^b$ (recall $a=r+1$). Pure powers span $\Sym^bV$ by
Lemma~\ref{lem:span}, and $T$ is multilinear, so $T=0$. This completes the
induction.
\end{proof}

\section{Proof of the theorem and a corollary}\label{sec:comparison}

\begin{proof}[Proof of Theorem~\ref{thm:main}]
Let $a\ge2$ and $b\ge a(a-1)$; in particular $b\ge1$. By
Proposition~\ref{prop:vanishing}, every symmetric $a$-linear form on
$\Sym^bV$ whose diagonal vanishes on all products of $b$ vectors is zero. By
Proposition~\ref{prop:criterion}, $\mu_{a,b,V}$ is surjective.
\end{proof}

The ordinary Foulkes conjecture asks only for the existence of an
equivariant injection in the opposite direction (see
Section~\ref{sec:context}). Surjectivity gives this in the range of the
theorem.

\begin{corollary}\label{cor:comparison}
For every finite-dimensional complex vector space $V$, every integer
$a\ge2$, and every integer $b\ge a(a-1)$, there exists a
$\mathrm{GL}(V)$-equivariant embedding
\[
 \Sym^a(\Sym^b V)\hookrightarrow\Sym^b(\Sym^a V).
\]
In particular, this gives the sixth symmetric-power comparison for $b\ge30$.
\end{corollary}

\begin{proof}
The group $\mathrm{GL}(V)$ acts on $\Sym V$ by algebra automorphisms and on
$\cP$ diagonally, commuting with $\Sa$. Averaging and multiplication
therefore commute with the action, so the surjection $\mu_{a,b,V}$ of
Theorem~\ref{thm:main} is equivariant. Finite-dimensional polynomial
representations of $\mathrm{GL}(V)$ over $\C$ are completely reducible
\cite[Chapter~I, Appendix~A, Sections~7--8]{Macdonald}. Choose an invariant
complement $M$ to $\ker\mu_{a,b,V}$. The restriction $\mu_{a,b,V}|_M$ is an
equivariant isomorphism onto $\Sym^a(\Sym^b V)$. Its inverse followed by the
inclusion $M\subseteq\Sym^b(\Sym^a V)$ is the required embedding. For
$a=6$ the range is $b\ge30$. This construction asserts existence only; it
does not identify the embedding with a prescribed canonical map.
\end{proof}

\begin{remark}[Scope]
Lemma~\ref{lem:shift} holds for any number $n\ge1$ of coordinates and
arbitrary additional variables, and the degree comparisons in
Proposition~\ref{prop:vanishing} depend only on $a$ and $b$. Together with
the finite-dimensional annihilator argument of
Proposition~\ref{prop:criterion}, this gives the same bound in every finite
dimension, without a stable-dimension reduction. The proof uses positivity
of the Fischer Hermitian inner product, the splitting of binary forms over
$\C$, and the fact that the scalar factors in the coefficient and
differentiation identities ($a$, $k!$, $b-k$, $b!/m!$) are nonzero in
characteristic zero.
\end{remark}

\section{Context and related work}\label{sec:context}

The statements in this section are background taken from the source,
with its citations and locators. They are not used in the proofs above, and
they were not checked against the cited works for this rewrite.

\paragraph{Formulation.}
In positive dimension, Raicu--Sam--Weyman \cite[Section~2]{RSW} formulate
the map as $\Sym^b(D^aV)\longrightarrow D^a(\Sym^bV)$, with divided powers
realized as invariant tensors. Transporting both sides by averaging gives
exactly \eqref{eq:map}.

\paragraph{The stabilization question.}
For fixed $a$ and $V$, one asks for a degree $b_0\ge0$ such that
$\mu_{a,b,V}$ is surjective for every $b\ge b_0$, either the least such
degree in a fixed dimension or a bound valid in every finite dimension. The
affine Chow variety of decomposable $a$-forms consists of products of $a$
linear forms in $\Sym^a(V^*)$. Its coordinate ring identifies with $A$,
whose normalization is $R$; here $b$ is the coordinate degree
\cite[Section~2, Theorem~2.3]{RSW}. See also Landsberg
\cite[Section~7.6, Propositions~7.11--7.12]{Landsberg}. This also explains
the reversed roles of the parameters: $T(P,\ldots,P)$ is a homogeneous
degree-$a$ polynomial on $\Sym^bV$, tested on products of $b$ vectors viewed
as linear forms on $V^*$; compare the product parametrization in
\cite[Section~2]{RSW}. Brion established eventual surjectivity
\cite{Brion1993} and subsequently gave an effective bound involving both $a$
and the dimension \cite[Theorem~3.3]{Brion}. Eventual surjectivity says that
the graded quotient $R/A$ has only finitely many nonzero components; the
least stabilization degree is one more than its highest nonzero degree, or
zero if $R=A$. Raicu, Sam, and Weyman bound the regularity of the Segre
ring and modules of covariants, including the normalization
\cite[Theorem~3.9]{RSW}; their theorem does not compute the top degree of
$R/A$. Theorem~\ref{thm:main} gives the sufficient bound $a(a-1)$ in every
dimension. It answers positively the request for a polynomial bound in
\cite[Problem~7.19]{Landsberg}: in that source's notation $d=b$, $n=a$,
$w=\dim V$, the bound is $d\ge n(n-1)$, including when $w=n$. The argument
does not determine the least stabilization degree, and the bound is not
asserted to be sharp.

\paragraph{Relation to the Foulkes conjecture.}
The ordinary Foulkes conjecture \cite{Foulkes1950} asks, for $a\le b$, for
a $\mathrm{GL}(V)$-equivariant injection
$\Sym^a(\Sym^b V)\hookrightarrow\Sym^b(\Sym^a V)$
\cite[Conjecture~2.10]{RSW}. Corollary~\ref{cor:comparison} gives this in
the range of Theorem~\ref{thm:main}. Existence of an embedding does not
specify the canonical map. At $a=b$ the two representations coincide, so the
comparison is automatic, while the canonical endomorphism can fail to be
invertible: it has nonzero kernels at $(a,b)=(5,5)$ and $(6,6)$ in suitable
dimensions
\cite[arXiv version, Theorem~6(b),(c) and Remark~7]{CIM}, both outside the
range of Theorem~\ref{thm:main}.

\paragraph{Related techniques.}
Successive weight shifts also occur in Ikenmeyer's proof
\cite[Section~5]{Ikenmeyer} of McKay's propagation theorem \cite{McKay};
his argument \cite[Section~5, Claim~5.1]{Ikenmeyer} decomposes tensors
according to untouched positions before applying a two-coordinate weight
shift. In Ikenmeyer's formulation, the canonical maps in the opposite
direction $\Sym^a(\Sym^b V)\longrightarrow\Sym^b(\Sym^a V)$ are considered
with $\dim V\ge b$ for each $b$; for fixed $a$ and $B\ge a$, injectivity at
$(a,B)$ propagates to every $(a,b)$ with $b\ge B$
\cite[Theorem~1.3 and Section~3]{Ikenmeyer}. Here the two derivative chains
(Steps~2 and~4) surround an induction on the number of slots and give a
surjectivity range directly.

\section{Verification checklist}

The steps that carry the weight of the proof are listed below, in the order
in which a sceptical reader should check them.
\begin{enumerate}[leftmargin=*]
\item \textbf{Lemma~\ref{lem:shift}, Steps 3--5.} The adjoint of $\D YU$ for
the Fischer inner product is $\D UY$; the commutator is
$D'D-DD'=N_U-N_Y$ (check the sign); and the norm identity \eqref{eq:norm}
follows. Everything else rests on this lemma.
\item \textbf{Proposition~\ref{prop:vanishing}, Step 1.} $(x^r+ct^r)Q$ is a
product of $b$ vectors for \emph{every} $c\in\C$ (this uses roots in $\C$),
and the coefficient of $c$ is $a\,T(x^rQ,\ldots,x^rQ,t^rQ)$.
\item \textbf{Step 2, items (2b)--(2c).} Before the operator
$\D{y_{k+1}}{t}$, the $t$-degree is $b-k$ and the $y_{k+1}$-degree is still
$r$ (not $r+1$), so $p-q=m-k\ge1$. The chain ends at $k=m$, where
$b-m=r$ matches \eqref{eq:binary}.
\item \textbf{Step 3.} $S_{x,t}$ is a symmetric $r$-linear form on
$\Sym^mV$, and the induction hypothesis is applied with parameters $(r,m)$,
which requires $m\ge1$ and $m\ge r(r-1)$; both follow from
\eqref{eq:range}.
\item \textbf{Step 4, item (4d).} All $r^2$ operators raise the degree of
the same block $x$, so $q\le r^2-1$, while $p\ge m+1\ge r^2+1$. This is the
place where $b\ge a(a-1)$ (that is, $m\ge r^2$) is used in full. Also check
in (4b) that $v_i$ occurs in one slot only, so that Lemma~\ref{lem:rule}
applies.
\item \textbf{Proposition~\ref{prop:criterion}, Steps B and C.} $A_b$ is
spanned by the tensors $P^{\otimes a}$ (this uses Lemma~\ref{lem:span} with
$k=a$), and $\lambda\mapsto T_\lambda$ is a bijection onto symmetric
$a$-linear forms.
\item \textbf{Lemma~\ref{lem:span}.} The sign count in the polarization
identity \eqref{eq:polarization}; it is used three times (criterion, base
case, final line).
\end{enumerate}

\begin{thebibliography}{99}
\bibitem{AldazRender}
J.~M.~Aldaz and H.~Render,
\href{https://doi.org/10.1007/s13324-023-00844-4}{\emph{A Fischer type decomposition theorem from the apolar inner product}},
Analysis and Mathematical Physics \textbf{13} (2023), article~91.

\bibitem{Brion1993}
M.~Brion,
\href{https://doi.org/10.1007/BF03026558}{\emph{Stable properties of plethysm: on two conjectures of Foulkes}},
Manuscripta Mathematica \textbf{80} (1993), no.~4, 347--371.

\bibitem{Brion}
M.~Brion,
\href{https://smf.emath.fr/publications/sur-certains-modules-gradues-associes-aux-produits-symetriques}{\emph{Sur certains modules gradu\'es associ\'es aux produits sym\'etriques}},
in \emph{Alg\`ebre non commutative, groupes quantiques et invariants}
(Reims, 1995), S\'eminaires et Congr\`es \textbf{2},
Soci\'et\'e Math\'ematique de France, Paris, 1997, 157--183.

\bibitem{CIM}
M.-W.~Cheung, C.~Ikenmeyer, and S.~Mkrtchyan,
\href{https://doi.org/10.1016/j.jsc.2016.09.002}{\emph{Symmetrizing tableaux and the 5th case of the Foulkes conjecture}},
Journal of Symbolic Computation \textbf{80} (2017), part~3, 833--843.
\href{https://arxiv.org/abs/1509.03944v1}{arXiv:1509.03944v1}.

\bibitem{Foulkes1950}
H.~O.~Foulkes,
\href{https://doi.org/10.1112/jlms/s1-25.3.205}{\emph{Concomitants of the quintic and sextic up to degree four in the coefficients of the ground form}},
Journal of the London Mathematical Society \textbf{25} (1950), no.~3,
205--209.

\bibitem{FultonHarris}
W.~Fulton and J.~Harris,
\href{https://doi.org/10.1007/978-1-4612-0979-9}{\emph{Representation Theory: A First Course}},
Graduate Texts in Mathematics \textbf{129}, Springer-Verlag, New York, 1991.

\bibitem{Ikenmeyer}
C.~Ikenmeyer,
\emph{On McKay's propagation theorem for the Foulkes conjecture},
\href{https://arxiv.org/abs/1509.04957v1}{arXiv:1509.04957v1}, 2015.

\bibitem{Landsberg}
J.~M.~Landsberg,
\href{https://doi.org/10.1007/s11565-014-0202-7}{\emph{Geometric complexity theory: an introduction for geometers}},
Annali dell'Universit\`a di Ferrara \textbf{61} (2015), no.~1, 65--117.
\href{https://arxiv.org/abs/1305.7387}{arXiv:1305.7387}.
Locators refer to the published version.

\bibitem{Macdonald}
I.~G.~Macdonald,
\href{https://doi.org/10.1093/oso/9780198534891.001.0001}{\emph{Symmetric Functions and Hall Polynomials}},
second edition, Oxford Mathematical Monographs, Clarendon Press, Oxford, 1995.

\bibitem{McKay}
T.~McKay,
\href{https://doi.org/10.1016/j.jalgebra.2007.12.003}{\emph{On plethysm conjectures of Stanley and Foulkes}},
Journal of Algebra \textbf{319} (2008), no.~5, 2050--2071.

\bibitem{RSW}
C.~Raicu, S.~V.~Sam, and J.~Weyman,
\href{https://doi.org/10.1007/s10013-021-00527-2}{\emph{On some modules supported in the Chow variety}},
Vietnam Journal of Mathematics \textbf{50} (2022), no.~2, 501--521.
\href{https://arxiv.org/abs/2108.10910v1}{arXiv:2108.10910v1}.
Locators refer to this preprint version.
\end{thebibliography}

\end{document}
